% compiled with latexmk -pdf main.tex
% needs: graphicx, booktabs, pgfplots, siunitx, hyperref, cleveref
% and the macros generated by inject_numbers.py
\documentclass[10pt,a4paper]{article}
\usepackage[margin=1in]{geometry}
\usepackage{graphicx}
\usepackage{booktabs}
\usepackage{pgfplots}
\pgfplotsset{compat=1.18}
\usepackage{siunitx}
\usepackage{hyperref}
\usepackage{cleveref}
\usepackage{amsmath,amssymb}
\usepackage{authblk}

% --- macros filled by inject_numbers.py ---
\newcommand{\runId}{??}
\newcommand{\baseModel}{??}
\newcommand{\nSeeds}{??}
\newcommand{\nTest}{??}
\newcommand{\nFamilies}{??}
\newcommand{\gatesTotal}{??}
\newcommand{\gatesPassed}{??}
\newcommand{\promotionStatus}{??}
\newcommand{\accFamilyMean}{??}
\newcommand{\accFamilySd}{??}
\newcommand{\accRandomMean}{??}
\newcommand{\splitGap}{??}
\newcommand{\refusalBase}{??}
\newcommand{\refusalMean}{??}
\newcommand{\refusalDelta}{??}
\newcommand{\eceMean}{??}
\newcommand{\overlapThirteen}{??}

% --- local helpers ---
\DeclareMathOperator{\E}{E}
\DeclareMathOperator{\Var}{Var}
\newcommand{\abstain}{\varnothing}

\title{Governed Type-Safe Triage:\\
       LoRA Distillation Under a Refusal-Preserving Promotion Gate}
\author{Stephen P. Lutar}
\affil{SZL Holdings \textbullet\ stephenlutar2@gmail.com \textbullet\ ORCID 0009-0001-0110-4173}
\date{Preprint \textbullet\ 22 September 2026}

\hypersetup{pdftitle={Governed Type-Safe Triage (preprint draft)}}

\begin{document}
\maketitle

\begin{abstract}
We fine-tuned a Qwen3.5-0.8B base model into a type-safe technical triage classifier
across five seeds, then blocked our own release. The adapters learned the task and
abstained less than the frozen base. A twelve-check promotion gate caught it: one
check compares the candidate's abstention behaviour to the base on the same inputs
and cannot be traded against accuracy. The run passes 11 of 12 checks and ships
with promotion status \promotionStatus{} at \gatesPassed{}/\gatesTotal{}. We also
report what a conventional random split would have claimed: accuracy
\accRandomMean{} vs. \accFamilyMean{} (\splitGap{} inflation) when whole template
families are held out. The contribution is the mechanism and the disclosed negative
verdict, not the adapter. Promotion status is derived from a signed run receipt,
so it cannot be asserted by the process that wanted to pass.
\end{abstract}

\section{The gap this closes}
A model that finished training is not a model that passed. Most published results
collapse that distinction: an author picks a checkpoint, measures it on a held-out
split, and reports the number. Three decisions stay invisible to the reader.
Which checkpoint, and on what basis. How the held-out split was built. What the
model gave up to get what it gained.

This paper makes those three decisions mechanical for one narrow task.
\emph{Type-safe triage} assigns an incoming technical report to a typed category
with a schema-valid payload, and routes to a REVIEW sink when the type cannot be
established. The task is small enough for a sub-billion-parameter adapter and
hostile enough to expose the failure we care about: a model rewarded for confident
classification learns to stop declining.

Each candidate adapter is submitted to a twelve-check gate. Every check is a
pass/fail predicate over a receipt the run wrote itself: library versions, GPU
capability, split identity, hyperparameters, contamination verdict, and artifact
digests. Release status is a function of that receipt. One check compares the
candidate's abstention behaviour to the frozen base on the same inputs, and it is
not tradeable against accuracy. That check is the one this run failed.

\textbf{What we are claiming.} Not that the adapter is good. The gate says it is
not cleared. The claim is narrower and harder: a pipeline whose release status is
derived from a signed receipt produces usable evidence even when the evidence is
against itself. We publish the artifact, the receipt, and the failing check
together so a reader can check the verdict rather than take it.

\textbf{Contributions.}
\begin{itemize}
  \item \textbf{Refusal-preserving promotion} (\cref{sec:gate}): abstention as a
        non-compensatory constraint; an unproven contamination check defaults to
        refusal.
  \item \textbf{A twelve-check gate specification} (\cref{sec:gate}): each check a
        predicate over a receipt, so release cannot be announced, only computed.
  \item \textbf{A family-held-out protocol and the cost of not using it}
        (\cref{sec:setup,sec:results}): measured inflation from random splitting on
        the same data pool.
  \item \textbf{An operations contract} (\cref{sec:ops}): burn-rate objectives for a
        deployed agentic triage pipeline, including the two guards that keep the
        alerts survivable.
  \item \textbf{A disclosed negative verdict}: the run at 11/12, published with the
        failing check named.
\end{itemize}

\section{Related work}
\textbf{Parameter-efficient adaptation.}
Low-rank adaptation (Hu et al., 2022) and its quantised form (Dettmers et al.,
2023) cut the cost of specialising a base model far enough that evaluation
discipline, not compute, became the bottleneck in small-model work. Our
contribution is not the adapter. It is what stands between the adapter and a
release.

\textbf{Contamination and split construction.}
Benchmark contamination is well documented in large-model evaluation (Dodge et
al., 2021; Sainz et al., 2023). Its small-data analogue is near-duplicate
leakage across a random split. Holding back groups instead of rows is standard
in domains with structured redundancy; we use template families as the group and
report the resulting gap rather than asserting the protocol is better.

\textbf{Refusal, abstention, calibration.}
Selective prediction (El-Yaniv and Wiener, 2010) frames abstention as a
coverage-risk trade-off. Alignment work treats refusal as a behaviour training
can remove (Bai et al., 2022). Calibration error (Guo et al., 2017) asks whether
stated confidence tracks accuracy. Our gate uses a refusal-preservation predicate
and a calibration bound separately, because they fail for different reasons.

\textbf{Provenance.}
in-toto (Torres-Arias et al., 2019) and SLSA describe verifiable build
provenance, with DSSE as the signing envelope. We apply the same posture to a
training run: the receipt is the record, and the model card may not contradict it.

\textbf{Operating trained models.}
Multi-window burn-rate alerting (Beyer et al., 2018) is the operational standard
for separating a real regression from a transient spike. \cref{sec:ops} applies it
where the failure modes are tool-call loops and runaway token spend rather than
HTTP errors.

\section{Formalisation}
Let $\mathcal{X}$ be incoming reports and $\mathcal{Y} = \{y_1,\dots,y_K\}
\cup \{\varnothing\}$ the typed label space, with $\varnothing$ the REVIEW sink.
A triage policy is a map $f : \mathcal{X} \to \mathcal{Y} \times \Sigma$ into labels
and schema-valid payloads.

\begin{definition}[Type safety]
$f$ is type-safe on $x$ if either $f(x) = (\varnothing,\emptyset)$, or
$f(x) = (y,\sigma)$ with $\sigma$ validating against the schema declared for $y$.
\end{definition}
Type safety constrains the output channel, not correctness. A type-safe policy
can be wrong; it cannot be unparseable. We enforce it by validation outside the
weights, so it holds as an invariant rather than a learned tendency.

\begin{definition}[Abstention profile]
For policy $f$ and set $D$, let $\rho(f,D)$ be the fraction of $D$ on which $f$
returns $\varnothing$. Let $\rho_S(f,D)$ restrict that average to a curated subset
$S \subseteq D$ of inputs that should be declined.
\end{definition}

\begin{definition}[Refusal-preserving promotion]
Let $f_0$ be the frozen base, $f_\theta$ a candidate, and $\tau \ge 0$ a tolerance.
The candidate preserves refusal iff
\[
\rho_S(f_\theta,D) \ge \rho_S(f_0,D) - \tau.
\]
Release requires this \emph{and} every other check. No gain in accuracy,
macro-$F_1$, or calibration may be exchanged for a violation.
\end{definition}

\begin{definition}[Derived promotion]
Let $r$ be a receipt, $G = \{g_1,\dots,g_{12}\}$ predicates over $r$, and
$v(r)$ the contamination verdict recorded inside $r$. Then
\[
\mathsf{status}(r) =
\begin{cases}
\mathtt{PROMOTABLE} & \text{if } \forall i\, g_i(r) \land v(r)=\mathtt{CLEAN}\\
\mathtt{NOTPROMOTABLE} & \text{otherwise}
\end{cases}
\]
\end{definition}

\begin{proposition}[Unproven defaults to refusal]
If $v(r)$ is absent or $\mathtt{UNPROVEN}$, then
$\mathsf{status}(r) = \mathtt{NOTPROMOTABLE}$.
\end{proposition}
\emph{Proof.} The first case requires $v(r)=\mathtt{CLEAN}$. Any other value,
including absence, falls to the second. $\square$

The operational content: a contamination check that never ran is treated exactly
like one that failed. Teams discover this rule the first time a pipeline crashes
before the check and the run still reports a pass.

\section{Method}
\subsection{Distillation}
The base $f_0$ is Qwen3.5-0.8B, weights frozen. Candidates are low-rank adapters
of \baseModel{} bytes, trained with loss on assistant turns only: the mask excludes
the prompt, so the model is scored on what it must produce rather than on text
handed to it. Rank and scaling are pinned to the predecessor run so cross-run
comparison means something.

\subsection{Seeds as a reported axis}
We train five adapters at seeds 11, 23, 37, 53, 71 and report seed spread
instead of averaging it away. A check that passes on three seeds and fails on
two is a different finding from one that fails everywhere, and a mean erases the
difference. \cref{fig:gate-ledger} keeps the per-seed resolution.

\subsection{Pipeline}
The order matters: the receipt is written before the gate runs, and the gate reads
only the receipt. Nothing downstream can observe a value that was never recorded.

\begin{verbatim}
for each seed s:
    theta_s <- TrainAdapter(f_0, s)   # loss on assistant turns
    e_s     <- Evaluate(theta_s, D_test)  # frozen protocol
    v_s     <- ContaminationCheck(D_train, D_test)
    r_s     <- WriteReceipt(theta_s, e_s, v_s, env, hashes)
    Sign(r_s)                         # DSSE, or UNSIGNED_HONEST
r <- Aggregate({r_s})
for each g_i in G:  p_i <- g_i(r)   # 12 predicates
return status(r), naming every i where p_i is false
\end{verbatim}

\subsection{Validation outside the weights}
Schema validation runs on the decoded output, and a failure becomes $\varnothing$
rather than a retry loop that edits the output into compliance. The model is never
credited for a payload a validator repaired. This costs accuracy and is why REVIEW
carries real support in \cref{fig:per-class}.

\section{The promotion gate}
Each predicate reads the receipt, each is reproducible from the archived bundle,
and each is reported per seed.

\begin{table}[h]
\centering\small
\caption{The twelve-check gate. Release needs the conjunction; the contamination
verdict is a separate precondition by Definition 4.}
\begin{tabular}{@{}clp{7cm}@{}}
\toprule
\# & Check & Predicate over the receipt \\
\midrule
1 & schema conformance & every payload validates, or the output is $\varnothing$ \\
2 & fixture replay & stored fixtures reproduce byte-identical decisions \\
3 & family-split integrity & no template family appears on both sides \\
4 & n-gram leakage bound & train/test overlap under threshold for $n=5,8,13,21$ \\
5 & refusal preservation & Definition 3 holds at tolerance $\tau$ \\
6 & abstain routing & unresolvable types reach $\varnothing$, never a default class \\
7 & per-class floor & no class below its recall floor \\
8 & calibration bound & expected calibration error under threshold \\
9 & latency budget & p95 within the serving budget \\
10 & token budget & tokens per decision within budget \\
11 & receipt signature & DSSE verifies, or state is UNSIGNED_HONEST \\
12 & weight hash pin & adapter SHA-256 matches the published manifest \\
\bottomrule
\end{tabular}
\label{tab:gate}
\end{table}

\textbf{Twelve checks, not one score.}
A weighted score over these quantities would be easier to optimise and would
destroy the property we want. Checks 5 and 6 exist to make one specific trade
impossible; folding them into an aggregate re-enables it. The gate blocks often.
That is the design working, not a defect to tune out.

\section{Receipts}
A receipt is JSON the run writes: environment (library versions, GPU capability
and architecture list), split identity and hashes, hyperparameters, the evaluation
block, the contamination verdict, and artifact digests. Receipts are hash-chained
and DSSE-signed where a key exists. Where no key exists the state is recorded as
UNSIGNED_HONEST rather than omitted, because an omitted signature reads as a
signed one to anyone skimming.

Two properties carry the paper. Release status is computed from the receipt, so a
console message claiming a pass has no standing. And the archived bundle carries
SHA-256 for every file, so a reader can confirm the files described are the files
provided.

\textbf{The audit that found our own error.}
An earlier transcript of this pipeline reported a passing 12/12 gate. The receipt
disagreed: blocked at 11/12. The message was a print statement, not a derived
value, and it survived because nobody had required it to be derived. Definition 4
is the fix. This paper reports the receipt.

\section{Setup}
\label{sec:setup}
\textbf{Data and splits.}
The evaluation set holds \nTest{} items across \nFamilies{} template families.
Partitioning is by family: items sharing a family stay on one side of the
boundary. We also build a conventional random split over the same pool, used only
to measure inflation in \cref{sec:results}, never for release.

\textbf{Contamination.}
We compute train/test n-gram overlap for $n=5,8,13,21$ and pairwise Jaccard
similarity across the family boundary. A CLEAN verdict requires overlap under
threshold at every $n$. Anything else records UNPROVEN.

\textbf{Frozen protocol.}
Decoding parameters and the schema set are frozen before training, and the
protocol hash goes into every receipt. Changing the protocol invalidates
cross-run comparison and counts as a new experiment.

\textbf{Hardware.}
Recorded per receipt. We record the GPU architecture list rather than an
availability flag, because availability can report true on an architecture whose
kernels were never compiled into the installed wheel. That distinction cost us a
full evaluation run.

\section{Results}
\label{sec:results}

\begin{figure}[htbp]
\centering
\includegraphics[width=\linewidth]{figures/fig1_loss.pdf}
\caption{Distillation loss across five seeds, loss masked to assistant turns.
The band spans the seed envelope. Synthetic placeholder.}
\label{fig:loss}
\end{figure}

\textbf{Verdict.}
The run passed 11 of 12 checks: status \promotionStatus{} at \gatesPassed{}/\gatesTotal{}.
The failing check is refusal preservation. \cref{fig:gate-ledger} gives the ledger
per seed and \cref{fig:refusal} the magnitude against the frozen base at
\refusalBase{}.

\begin{figure}[htbp]
\centering
\includegraphics[width=\linewidth]{figures/fig2_gate_ledger.pdf}
\caption{Gate ledger. Green is PASS, red is FAIL. Per-seed resolution is kept
deliberately: a check that fails on a subset is a different problem from one that
fails everywhere. Synthetic placeholder.}
\label{fig:gate-ledger}
\end{figure}

\begin{figure}[htbp]
\centering
\includegraphics[width=\linewidth]{figures/fig3_refusal.pdf}
\caption{Refusal preservation. Bars give candidate minus base refusal rate;
negative bars violate Definition 3. Candidate mean \refusalMean{} against base
\refusalBase{}, a change of \refusalDelta{}.
}
\label{fig:refusal}
\end{figure}

\textbf{What random splitting would have told us.}
Family-split accuracy is \accFamilyMean{} ($\pm$\accFamilySd{} across seeds).
Random-split accuracy over the same pool is \accRandomMean{}: a gap of
\splitGap{}. Every seed moves the same direction, which is what leakage looks
like and not what sampling noise looks like.

\begin{figure}[htbp]
\centering
\includegraphics[width=0.85\linewidth]{figures/fig4_split_gap.pdf}
\caption{Paired split comparison. Each line is one seed. Synthetic placeholder.}
\label{fig:split-gap}
\end{figure}

\begin{figure}[htbp]
\centering
\includegraphics[width=0.65\linewidth]{figures/fig5_calibration.pdf}
\caption{Reliability diagram, mean expected calibration error \eceMean{}.
}
\label{fig:calibration}
\end{figure}

\begin{figure}[htbp]
\centering
\includegraphics[width=\linewidth]{figures/fig6_contamination.pdf}
\caption{Contamination. Left: train/test n-gram overlap, with 13-gram overlap at
\overlapThirteen{}. Right: pairwise Jaccard similarity across the family
boundary. Synthetic placeholder.}
\label{fig:contamination}
\end{figure}

\begin{figure}[htbp]
\centering
\includegraphics[width=\linewidth]{figures/fig7_per_class.pdf}
\caption{Per-class precision and recall with support. REVIEW is an outcome, not
an error bucket. Synthetic placeholder.}
\label{fig:per-class}
\end{figure}

\begin{figure}[htbp]
\centering
\includegraphics[width=\linewidth]{figures/fig8_serving.pdf}
\caption{Serving cost: latency percentiles and token budget per evaluation pass,
the two quantities monitored as error budgets in \cref{sec:ops}. Synthetic
placeholder.}
\label{fig:serving}
\end{figure}

\textbf{Reading the negative result.}
The adapters learned the task. Loss fell consistently and family-split accuracy
beat the base. They also declined less often. Under Definition 3 that
disqualifies them, and no number in this section changes the verdict.

\section{Sensitivity}
\textbf{Tolerance.}
$\tau$ sets how much abstention loss is acceptable. We report the verdict as a
function of $\tau$ rather than defending one value: this run clears check 5 only
for $\tau$ above the largest per-seed regression in \cref{fig:refusal}. A reader
who prefers a looser threshold can locate their own answer, which is the honest
form of a threshold claim.

\textbf{Seed count.}
With five seeds, the ledger separates systematic failure from incidental. A
single-seed run reports a pass or a fail with no basis for telling those apart,
which is most published fine-tuning.

\textbf{Loss masking.}
Masking to assistant turns reduces effective tokens per example. We report the
resulting token budget in \cref{fig:serving} rather than filing it under
implementation detail.

\section{Operating the model}
\label{sec:ops}
A promoted policy runs inside an agentic pipeline, where the failures are
tool-call loops, one degraded tool, and runaway token spend, not HTTP 500s. We
monitor error-budget burn instead of thresholds.

Let the objective be that a fraction $1-\epsilon$ of runs succeed over a
30-day window. For window $W$, let $\beta_W$ be the observed bad-run ratio; the
burn rate is $\beta_W / \epsilon$. An alert fires only when a long window and its
short confirmation window both exceed the threshold.

\begin{table}[h]
\centering\small
\caption{Burn-rate tiers at $\epsilon=0.01$. The short window is one twelfth of
the long window.}
\begin{tabular}{@{}lcccl@{}}
\toprule
Tier & Windows & Burn rate & Budget consumed & Route \\
\midrule
1 & 1 h + 5 min  & 14.4 & 2\% per hour & page \\
2 & 6 h + 30 min & 6.0  & 5\% in 6 h & page \\
3 & 1 d + 2 h    & 3.0  & 10\% in 1 d & ticket \\
4 & 3 d + 6 h    & 1.0  & trend drain & review \\
\bottomrule
\end{tabular}
\label{tab:burn}
\end{table}

Instrumentation follows the OpenTelemetry generative-AI conventions: one trace
per triage run, spans per tool and model call, token usage as a metric split by
type. Two guards keep the signal usable. A minimum-traffic floor, so one failure
in a quiet window cannot read as a total outage. And exclusion of evaluation
traffic, so benchmark sweeps never spend the production error budget.

\textbf{Why not alert on latency.}
A p95 threshold fires on spikes that resolve themselves and stays quiet during
slow decay that drains the budget. Express latency as a compliance ratio, apply
the same tiers, and the alert tracks budget consumption instead of weather.

\section{What a blocked release costs}
A gate that blocks has a bill, and pretending otherwise is how gates get
disabled six months after they are installed.

\textbf{Direct cost.}
This run produced five trained adapters, a full evaluation pass, and no shippable
artifact. The compute is spent. The engineering time is spent. What remains is a
measured artifact and a named reason it cannot ship.

\textbf{The alternative cost.}
Shipping it means deploying a triage classifier that declines less often than the
base model on inputs curated to require declining. In a triage system the
consequence is specific: reports that should have reached a human get a confident
type instead, and the failure surfaces as a misrouted incident rather than as an
error.

\textbf{When overriding is defensible.}
Three conditions together, not separately. The abstention subset is wrong for
your deployment and you can say how. The downstream consumer of REVIEW has
independent review capacity. The override is recorded in the receipt with a named
owner, so the next audit finds a decision rather than a silence. Absent all
three, the block stands.

\textbf{What to do next with a blocked artifact.}
Publish it, then fix the objective rather than the threshold. The regression here
is a training-signal problem: the loss rewards committing to a type and never
rewards declining. Adding abstention-positive examples to the mixture is a
hypothesis with a measurable prediction, and the gate is already set up to test
it.

\section{Threats to validity}
\textbf{Family assignment is our choice.}
We define template families. A coarser or finer definition moves the reported
gap. The assignment ships with the bundle so the choice is inspectable, but it is
not derived from first principles.

\textbf{The refusal subset is curated.}
$S$ is hand-built. A different curation gives a different abstention profile and
possibly a different verdict on check 5. This is the most consequential judgement
in the paper and the first thing a reviewer should attack.

\textbf{One base, one task.}
Nothing here shows that refusal regression under distillation generalises. It
happened here, and the gate caught it here.

\textbf{Signing state.}
Receipts recorded UNSIGNED_HONEST are tamper-evident by hash chain only, not by
signature.

\section{Limitations}
The pipeline is deliberately conservative. Unproven contamination counts as
contamination. Repaired payloads count as refusals. One failing check blocks the
release. Each trades headline performance for auditability, and anyone optimising
for a leaderboard should not adopt them. The gate is also not complete: twelve
checks are the twelve we could derive from a receipt, not an account of
everything that can go wrong.

\section{Reproducing this}
The archived bundle holds evaluation outputs, gate results, receipts, and a
manifest with SHA-256 per file. Figures are generated from that data by a script
that refuses to render a panel whose inputs are missing, so an absent measurement
appears as an absent figure rather than a plausible curve. The archival record is
versioned; corrections appear as new versions under the same concept identifier
rather than as edits.

To reproduce: clone the source, install the pinned environment, replay the
fixtures (check 2), recompute the contamination verdict (checks 3 and 4), and
re-evaluate the published adapters. If your receipt disagrees with ours, that is
our defect and the issue tracker is the right destination.

\section{Conclusion}
We specialised a small model for type-safe triage, measured it with whole
template families held out, and blocked the release because it declined less than
the model it came from. The artifact ships as NOTPROMOTABLE at 11/12 with the
failing check named and the receipt attached.

The useful part is not the adapter. It is that the verdict was computed rather
than announced, which is what made it possible for the pipeline to return an
answer its authors did not want.

\section*{Data and code availability}
Source, receipts, and gate ledger under Apache-2.0; text and figures under
CC-BY-4.0. The archived bundle and its citable record are referenced in the
accompanying manifest.

\bibliographystyle{plain}
\bibliography{refs}

\appendix
\section{Gate predicate specifications}
\label{app:gates}
\begin{itemize}
  \item \textbf{G01 schema conformance.} Reads eval.payload_validation. Fails if any
        payload is unparseable or violates its schema. Catches a model that emits
        confident prose where a typed object was required.
  \item \textbf{G02 fixture replay.} Reads fixtures.digest and replays the stored
        decision set. Fails on byte-level divergence. Catches nondeterminism from
        decoding settings or library drift.
  \item \textbf{G03 family-split integrity.} Reads split.family_map. Fails if a
        family appears on both sides. Catches leakage introduced by re-partitioning
        after the fact.
  \item \textbf{G04 n-gram leakage bound.} Reads contamination.ngram. Fails above
        threshold at any recorded $n$. Catches near-duplicates that survive family
        separation.
  \item \textbf{G05 refusal preservation.} Reads eval.refusal_rate for candidate and
        base. Fails on violation of Definition 3. Catches abstention trained away
        by an accuracy objective.
  \item \textbf{G06 abstain routing.} Reads eval.review_routing_rate and the
        unresolvable-type fixtures. Fails if an unresolvable input receives a
        typed label. Catches a default class absorbing uncertainty.
  \item \textbf{G07 per-class floor.} Reads eval.per_class. Fails if a class drops
        below its recall floor. Catches aggregate accuracy hiding a collapsed
        minority class.
  \item \textbf{G08 calibration bound.} Reads eval.ece and eval.confidence_bins.
        Catches confidence that does not track correctness.
  \item \textbf{G09 latency budget.} Reads eval.latency_ms_p95. Catches a quality
        gain bought with unshippable latency.
  \item \textbf{G10 token budget.} Reads eval.tokens_in and tokens_out. Catches
        context growth that quietly multiplies cost.
  \item \textbf{G11 receipt signature.} Reads the DSSE envelope; records
        UNSIGNED_HONEST when no key is present. Catches post-hoc receipt editing.
  \item \textbf{G12 weight hash pin.} Reads artifact.sha256. Catches a published
        adapter that is not the evaluated adapter.
\end{itemize}

\section{Receipt schema}
\label{app:receipt}
\begin{verbatim}
{
  "run_id": "...",
  "base_model": "Qwen3.5-0.8B",
  "adapter": {"seed": 11, "bytes": 25587104, "sha256": "..."},
  "env": {"torch": "...", "unsloth": "...", "gpu": "...", "arch_list": ["..."]},
  "split": {"protocol_hash": "...", "family_map": "...", "n_test": 0, "n_families": 0},
  "hyperparameters": {"r": 0, "alpha": 0, "lr": 0.0, "loss_mask": "assistant_only"},
  "contamination": {"verdict": "CLEAN|UNPROVEN|CONTAMINATED",
                     "ngram": [{"n": 13, "train_test_overlap_frac": 0.0}]},
  "eval": {"accuracy_family_split": 0.0, "accuracy_random_split": 0.0,
            "refusal_rate": 0.0, "review_routing_rate": 0.0, "ece": 0.0,
            "per_class": [{"label": "...", "precision": 0.0, "recall": 0.0, "support": 0}],
            "latency_ms_p50": 0.0, "latency_ms_p95": 0.0,
            "tokens_in": 0, "tokens_out": 0},
  "gates": [{"id": 5, "name": "refusal preservation", "status": "FAIL", "detail": "..."}],
  "promotion_status": "NOTPROMOTABLE",
  "signature": {"type": "DSSE", "state": "SIGNED|UNSIGNED_HONEST", "chain_prev": "..."}
}
\end{verbatim}

\section{Per-seed measurements}
\label{app:per-seed}
\input{results/per_seed_table}

\section{Hyperparameters and environment}
Injected from the receipt at build time. Rank and scaling are pinned to the
predecessor run. Learning rate, schedule, batch composition, sequence length,
GPU model, architecture list, and library versions reproduce verbatim from the
receipt env block, so a reader can rebuild the run without reading prose.

\section{Prompts and schemas}
The system prompt, the typed schema set, and one worked example per family,
including a family whose correct outcome is REVIEW.
\end{document}
